\section{问题三模型的建立与求解}

问题三在问题二容量硬约束方案已经确定的基础上研究服务定价。定价既影响老人价格满意度，也影响服务站收入、政府补贴和利润率。本文采用国赛优秀论文中常见的写法：先给出可复核的定价变量、补贴约束和利润率约束，再给出枚举求解过程，最后用分站财务表检验方案是否满足保本微利要求。

\subsection{模型的建立}

\subsubsection{定价变量}

基准选址方案为
\begin{equation}
J_0=\{C,E,G,I,J\},
\end{equation}
对应规模为小型、小型、小型、小型和大型。设可定价服务集合为
\begin{equation}
\mathcal{M}_p=\{\text{助餐},\text{日间照料},\text{上门护理},\text{康复理疗},\text{助浴}\}.
\end{equation}
紧急救助为公益免费项目，价格固定为 0，不参与定价优化，但其直接服务支出仍进入成本核算。设第 $m$ 类收费服务的基准价格为 $p_m^0$，服务站 $j$ 对该服务的优化价格为 $p_{j,m}$。

\subsubsection{价格满意度函数}

根据附件满意度规则，价格满意度为分段函数：
\begin{equation}
S_{j,m}^p(p_{j,m})=
\begin{cases}
1, & p_{j,m}\leq p_m^0,\\
0.90, & p_m^0<p_{j,m}\leq 1.1p_m^0,\\
0.75, & 1.1p_m^0<p_{j,m}\leq 1.2p_m^0,\\
0.60, & p_{j,m}>1.2p_m^0.
\end{cases}
\end{equation}
由于同一价格满意度区间内继续提高价格不会提高老人满意度，最优价格只需要在各满意度区间边界及必要降价点上搜索。本文设置候选价格集合为
\begin{equation}
\mathcal{P}_m=\{0.6p_m^0,0.8p_m^0,p_m^0,1.1p_m^0,1.2p_m^0\}.
\end{equation}
这一候选集的作用是同时覆盖降价释放潜在需求、基准价格运营和小幅上浮维持利润三类情形。

\subsubsection{政府补贴与利润率约束}

设服务站 $j$ 对服务 $m$ 的年有效服务人次为 $E_{j,m}$，单次直接支出为 $v_m$，年度运营成本为 $O_j$。政府对除紧急救助外的服务按有效服务人次补贴 2 元，并设置单站日补贴上限。记服务站 $j$ 的日补贴上限为 $L_j$，则年补贴为
\begin{equation}
G_j=365\min\left\{2\sum_{m\in\mathcal{M}_p}E_{j,m}^{\mathrm{day}},L_j\right\}.
\end{equation}
服务站 $j$ 年收入为
\begin{equation}
R_j=\sum_{m\in\mathcal{M}_p}p_{j,m}E_{j,m}.
\end{equation}
年直接服务成本为
\begin{equation}
C_j^{\mathrm{dir}}=\sum_{m=1}^{6}v_mE_{j,m}.
\end{equation}
年补贴后利润为
\begin{equation}
\Pi_j=R_j+G_j-C_j^{\mathrm{dir}}-O_j.
\end{equation}
根据保本微利和利润率不超过 $8\%$ 的要求，约束为
\begin{equation}
0\leq \rho_j=\frac{\Pi_j}{O_j+C_j^{\mathrm{dir}}}\leq 8\%,\quad j\in J_0.
\end{equation}

\subsubsection{目标函数}

小区 $i$ 的价格满意度按其服务需求结构加权。设小区 $i$ 的主服务站为 $j(i)$，则
\begin{equation}
S_i^p=\frac{\sum_{m\in\mathcal{M}_p}D_{i,m}^{\mathrm{bd}}S_{j(i),m}^p(p_{j(i),m})}{\sum_{m\in\mathcal{M}_p}D_{i,m}^{\mathrm{bd}}}.
\end{equation}
综合满意度为
\begin{equation}
S_i=0.2S_i^d+0.3S_i^r+0.5S_i^p.
\end{equation}
定价模型以老人规模加权综合满意度最大为目标：
\begin{equation}
\max \bar{S}_p=\frac{\sum_{i=1}^{10}N_iS_i}{\sum_{i=1}^{10}N_i}.
\end{equation}

\subsection{模型的求解}

\subsubsection{临界点枚举算法}

对每一个服务站，五类收费服务各有 5 个候选价格，总组合数为
\begin{equation}
5^5=3125.
\end{equation}
本文对每个服务站分别枚举全部价格组合，依次计算价格满意度、潜在需求释放、容量约束后的实际服务量、服务收入、政府补贴、直接成本、运营成本和利润率。删除利润率不满足 $[0,8\%]$ 的组合后，在剩余组合中选择价格满意度最高的方案；若价格满意度相同，则选择利润率更接近但不超过 $8\%$ 的方案。该算法没有随机参数，便于复核。

\subsubsection{定价结果}

\begin{table}[H]
\centering
\caption{问题三基准情景各站点部分服务最优价格与价格满意度}
\label{tab:p3_price}
\begin{tabular}{c c r r c}
\toprule
站点 & 服务 & 基准价格 & 优化价格 & 价格满意度 \\
\midrule
C & 助餐 & 10.00 & 6.00 & 1.0000 \\
C & 日间照料 & 20.00 & 20.00 & 1.0000 \\
C & 上门护理 & 30.00 & 30.00 & 1.0000 \\
C & 康复理疗 & 28.00 & 28.00 & 1.0000 \\
C & 助浴 & 25.00 & 25.00 & 1.0000 \\
E & 助餐 & 10.00 & 6.00 & 1.0000 \\
G & 助餐 & 10.00 & 8.00 & 1.0000 \\
全体站点 & 紧急救助 & 0.00 & 0.00 & 1.0000 \\
\bottomrule
\end{tabular}
\end{table}

表 \ref{tab:p3_price} 显示，最优价格组合没有通过提高价格来换取更高利润，而是在政府补贴约束下尽量保持或降低老人支付价格，使价格满意度保持为 1.0000。该结果与嵌入式养老服务的公益属性一致，也说明补贴政策能够在一定程度上吸收服务站运营压力。

\begin{table}[H]
\centering
\caption{问题三基准情景服务站补贴后财务结果}
\label{tab:p3_finance}
\begin{tabular}{c c r r r r}
\toprule
站点 & 规模 & 年补贴/元 & 补贴后利润/元 & 利润率 & 可行性 \\
\midrule
C & 小型 & 365000 & 56403.21 & 7.6324\% & 可行 \\
E & 小型 & 365000 & 44181.95 & 5.9786\% & 可行 \\
G & 小型 & 365000 & 53713.41 & 7.2684\% & 可行 \\
I & 小型 & 365000 & 56447.12 & 7.6383\% & 可行 \\
J & 大型 & 949000 & 127474.90 & 7.8277\% & 可行 \\
\midrule
合计/平均 & -- & 2409000 & 338220.60 & 7.2691\% & 可行 \\
\bottomrule
\end{tabular}
\end{table}

由表 \ref{tab:p3_finance} 可知，基准方案下年度政府补贴总额为 2409000 元，补贴后年度利润合计为 338220.60 元，平均利润率为 7.2691\%，最低利润率为 5.9786\%，全部站点满足 $0\%$ 至 $8\%$ 的保本微利约束。与仅计算总利润相比，分站利润率检验更能证明方案不存在单站亏损或单站过度盈利问题。

\subsubsection{结果解释与可及性分析}

从价格结构看，各站点均保持价格满意度 1.0000，说明最优解优先保证老人支付可及性。C、E 两个小型站将助餐价格降至 6 元，主要作用是释放基础生活服务需求；G 站助餐价格为 8 元，在维持价格满意度的同时提高了保本能力。J 站作为大型站承担更多跨小区服务量，获得 949000 元年补贴，利润率为 7.8277\%，接近但未超过上限，说明大型站在容量兜底和财务平衡中具有关键作用。

对自理老人而言，助餐和日间照料需求占比较高，低价助餐直接提高日常生活服务可及性。对半失能老人而言，上门护理和康复理疗需求更高，维持基准价并叠加政府补贴有助于控制服务负担。对失能老人而言，助浴、上门护理和紧急救助需求更集中，其可及性不仅取决于价格，还取决于问题二中的站点容量和响应能力。因此，问题三的定价结果应与容量硬约束选址方案联合解释。